\documentclass[11pt]{article}
\usepackage[a4paper,margin=25mm]{geometry}
\usepackage{fontspec}
\usepackage[hidelinks]{hyperref}
\setmainfont{Latin Modern Roman}
\pagestyle{plain}
\begin{document}
{\LARGE\bfseries ALIUS Bulletin - Guidelines for reviews - 2022\par}
\vspace{1em}
\section*{Alius Bulletin}
\section*{*}
\section*{Guidelines For}
\section*{Reviews}
If you want to edit a review for the ALIUS Bulletin, please read the following guidelines carefully.\par

Reviews offer an outlook on published works on consciousness and their relevance in respect to the diversity of consciousness. It is intended as short commentaries (500-1000 words), composed of the summarized main ideas of the work (250-500 words) and its relevance for the study of the diversity of consciousness (250-500 words).\par

Length of the commentary\par

Each Commentary (or reply to the commentary) must include the following:\par

-A title\par

-Main text: 500 < 1000 words\par

-References: no limit on the number of references; but the style used should be the 6th edition of the APA (if you use Zotero, this citation style can be downloaded at the following link: https://www.zotero.org/styles?q=id\%3Aapa).\par

-Addresses and affiliations of the authors for correspondence\par

-Font: Times New Roman; size: 12\par

Content of the commentary\par

Reviews are intended to highlight how theories of consciousness, including recently formulated ones, fit with the diversity of consciousness. It is intended as both highlighting some recent progress in consciousness science and providing a perspective about how they can inform the study of the diversity of consciousness. After having described one or several main ideas of the article/book, reviews can depart from the message of the original author and critically assess how the diversity of consciousness, or at least one specific conscious state, is given a new perspective under the light of this theory. It can also discuss unresolved theoretical points of the article/book from the perspective of the diversity of consciousness.\par

By doing so, reviews aim at raising the importance of embracing the diversity of consciousness when discussing the nature of cognition and/or consciousness. Each comment should be easily understandable to any scholar. For example, an anthropologist reading a comment conducted with a neuroscientist should easily understand most of the discussion, and conversely, a neuroscientist reading an interview conducted with an anthropologist should easily understand most of the discussion.\par

Modus operandi and requirements\par

1. Ask one of the editors of the Bulletin whether the article or book you want to review is deemed relevant as an open review.\par

2. Once editors have received the review, the editors will forward it to a native English-speaking proofreader who will take care of checking grammar and linguistic correctness. Note that the job of proofreaders is not to correct typos. As indicated above, typos must be corrected before sending the file to the editors.\par

3. The editors will next send you the feedback provided by the proofreaders. You will be expected to integrate the changes suggested by the proofreaders.\par

4. Once you have the author(s) agreement for publication, send the final version of the review to one of the editors. The review will then be ready to be published!\par

ALIUS contact email: aliusresearchgroup@gmail.com\par

Editorial team\par

Daniel Friedman (danielarifriedman@gmail.com); Matthieu Koroma (matthieu.koroma@gmail.com) ; George Fejer (gyorgyf@live.com)\par

\end{document}
